% Vorwärts- vs. inverse Kinematik am ebenen Arm mit drei Drehgelenken.
% Eigene Darstellung in Anlehnung an Kucuk2006 (Abb. 5a, S. 126: ebener Arm, Winkelkonvention;
% Abb. 1, S. 118: Vorwärts- vs. inverse Kinematik).
% Stil: ba-grafik.tex. Einbinden: % Vorwärts- vs. inverse Kinematik am ebenen Arm mit drei Drehgelenken.
% Eigene Darstellung in Anlehnung an Kucuk2006 (Abb. 5a, S. 126: ebener Arm, Winkelkonvention;
% Abb. 1, S. 118: Vorwärts- vs. inverse Kinematik).
% Stil: ba-grafik.tex. Einbinden: % Vorwärts- vs. inverse Kinematik am ebenen Arm mit drei Drehgelenken.
% Eigene Darstellung in Anlehnung an Kucuk2006 (Abb. 5a, S. 126: ebener Arm, Winkelkonvention;
% Abb. 1, S. 118: Vorwärts- vs. inverse Kinematik).
% Stil: ba-grafik.tex. Einbinden: % Vorwärts- vs. inverse Kinematik am ebenen Arm mit drei Drehgelenken.
% Eigene Darstellung in Anlehnung an Kucuk2006 (Abb. 5a, S. 126: ebener Arm, Winkelkonvention;
% Abb. 1, S. 118: Vorwärts- vs. inverse Kinematik).
% Stil: ba-grafik.tex. Einbinden: \input{Bilder/tikz/ik-vorwaerts-invers} in einer figure-Umgebung.
\begin{tikzpicture}[ba]
  % #1 = Stil der Gelenkwinkel, #2 = Zusatz an den Winkeln (z.B. "=?"), #3 = Stil Endeffektor-Label, #4 = Zusatz am Endeffektor
  \newcommand{\baIKArm}[4]{%
    \coordinate (B)  at (0,0);
    \coordinate (J1) at (20:2.3);
    \coordinate (J2) at ($(J1)+(50:2.0)$);
    \coordinate (E)  at ($(J2)+(78:1.8)$);
    % Bezugslinien: Horizontale an der Basis, Verlängerung der Glieder
    \coordinate (H)  at (2.4,0);
    \coordinate (X1) at ($(J1)+(20:1.2)$);
    \coordinate (X2) at ($(J2)+(50:1.1)$);
    \draw[ba hilfslinie, #1] (B) -- (H);
    \draw[ba hilfslinie, #1] (J1) -- (X1);
    \draw[ba hilfslinie, #1] (J2) -- (X2);
    % Glieder
    \draw[ba glied] (B) -- (J1) -- (J2) -- (E);
    % Gelenkwinkel
    \pic[ba winkel, #1, "$\theta_1#2$", angle eccentricity=1.12, pic text options={anchor=west}] {angle = H--B--J1};
    \pic[ba winkel, #1, "$\theta_2#2$", angle eccentricity=1.12, pic text options={anchor=west}] {angle = X1--J1--J2};
    \pic[ba winkel, #1, "$\theta_3#2$", angle eccentricity=1.12, pic text options={anchor=south west}] {angle = X2--J2--E};
    % Gelenke, Basis, Endeffektor
    \node[ba basis]  at (B)  {};
    \node[ba gelenk] at (J1) {};
    \node[ba gelenk] at (J2) {};
    \node[ba endeffektor] at (E) {};
    \node[right=1.5mm of E, #3] {$\vec{r}_\mathrm{E}#4$};
  }
  % a) Vorwärtskinematik: Winkel gegeben, Endeffektorlage gesucht
  \begin{scope}
    \baIKArm{}{}{ba gesucht}{=\,?}
    \node[ba teiltitel, align=center] at (2.0,-0.6) {a) Vorwärtskinematik};
  \end{scope}
  % Trennlinie
  \draw[ba trenner] (5.4,-1.0) -- (5.4,4.9);
  % b) Inverse Kinematik: Endeffektorlage gegeben, Winkel gesucht
  \begin{scope}[xshift=6.6cm]
    \baIKArm{ba gesucht}{=\,?}{}{}
    \node[ba teiltitel, align=center] at (2.0,-0.6) {b) Inverse Kinematik};
  \end{scope}
\end{tikzpicture}
 in einer figure-Umgebung.
\begin{tikzpicture}[ba]
  % #1 = Stil der Gelenkwinkel, #2 = Zusatz an den Winkeln (z.B. "=?"), #3 = Stil Endeffektor-Label, #4 = Zusatz am Endeffektor
  \newcommand{\baIKArm}[4]{%
    \coordinate (B)  at (0,0);
    \coordinate (J1) at (20:2.3);
    \coordinate (J2) at ($(J1)+(50:2.0)$);
    \coordinate (E)  at ($(J2)+(78:1.8)$);
    % Bezugslinien: Horizontale an der Basis, Verlängerung der Glieder
    \coordinate (H)  at (2.4,0);
    \coordinate (X1) at ($(J1)+(20:1.2)$);
    \coordinate (X2) at ($(J2)+(50:1.1)$);
    \draw[ba hilfslinie, #1] (B) -- (H);
    \draw[ba hilfslinie, #1] (J1) -- (X1);
    \draw[ba hilfslinie, #1] (J2) -- (X2);
    % Glieder
    \draw[ba glied] (B) -- (J1) -- (J2) -- (E);
    % Gelenkwinkel
    \pic[ba winkel, #1, "$\theta_1#2$", angle eccentricity=1.12, pic text options={anchor=west}] {angle = H--B--J1};
    \pic[ba winkel, #1, "$\theta_2#2$", angle eccentricity=1.12, pic text options={anchor=west}] {angle = X1--J1--J2};
    \pic[ba winkel, #1, "$\theta_3#2$", angle eccentricity=1.12, pic text options={anchor=south west}] {angle = X2--J2--E};
    % Gelenke, Basis, Endeffektor
    \node[ba basis]  at (B)  {};
    \node[ba gelenk] at (J1) {};
    \node[ba gelenk] at (J2) {};
    \node[ba endeffektor] at (E) {};
    \node[right=1.5mm of E, #3] {$\vec{r}_\mathrm{E}#4$};
  }
  % a) Vorwärtskinematik: Winkel gegeben, Endeffektorlage gesucht
  \begin{scope}
    \baIKArm{}{}{ba gesucht}{=\,?}
    \node[ba teiltitel, align=center] at (2.0,-0.6) {a) Vorwärtskinematik};
  \end{scope}
  % Trennlinie
  \draw[ba trenner] (5.4,-1.0) -- (5.4,4.9);
  % b) Inverse Kinematik: Endeffektorlage gegeben, Winkel gesucht
  \begin{scope}[xshift=6.6cm]
    \baIKArm{ba gesucht}{=\,?}{}{}
    \node[ba teiltitel, align=center] at (2.0,-0.6) {b) Inverse Kinematik};
  \end{scope}
\end{tikzpicture}
 in einer figure-Umgebung.
\begin{tikzpicture}[ba]
  % #1 = Stil der Gelenkwinkel, #2 = Zusatz an den Winkeln (z.B. "=?"), #3 = Stil Endeffektor-Label, #4 = Zusatz am Endeffektor
  \newcommand{\baIKArm}[4]{%
    \coordinate (B)  at (0,0);
    \coordinate (J1) at (20:2.3);
    \coordinate (J2) at ($(J1)+(50:2.0)$);
    \coordinate (E)  at ($(J2)+(78:1.8)$);
    % Bezugslinien: Horizontale an der Basis, Verlängerung der Glieder
    \coordinate (H)  at (2.4,0);
    \coordinate (X1) at ($(J1)+(20:1.2)$);
    \coordinate (X2) at ($(J2)+(50:1.1)$);
    \draw[ba hilfslinie, #1] (B) -- (H);
    \draw[ba hilfslinie, #1] (J1) -- (X1);
    \draw[ba hilfslinie, #1] (J2) -- (X2);
    % Glieder
    \draw[ba glied] (B) -- (J1) -- (J2) -- (E);
    % Gelenkwinkel
    \pic[ba winkel, #1, "$\theta_1#2$", angle eccentricity=1.12, pic text options={anchor=west}] {angle = H--B--J1};
    \pic[ba winkel, #1, "$\theta_2#2$", angle eccentricity=1.12, pic text options={anchor=west}] {angle = X1--J1--J2};
    \pic[ba winkel, #1, "$\theta_3#2$", angle eccentricity=1.12, pic text options={anchor=south west}] {angle = X2--J2--E};
    % Gelenke, Basis, Endeffektor
    \node[ba basis]  at (B)  {};
    \node[ba gelenk] at (J1) {};
    \node[ba gelenk] at (J2) {};
    \node[ba endeffektor] at (E) {};
    \node[right=1.5mm of E, #3] {$\vec{r}_\mathrm{E}#4$};
  }
  % a) Vorwärtskinematik: Winkel gegeben, Endeffektorlage gesucht
  \begin{scope}
    \baIKArm{}{}{ba gesucht}{=\,?}
    \node[ba teiltitel, align=center] at (2.0,-0.6) {a) Vorwärtskinematik};
  \end{scope}
  % Trennlinie
  \draw[ba trenner] (5.4,-1.0) -- (5.4,4.9);
  % b) Inverse Kinematik: Endeffektorlage gegeben, Winkel gesucht
  \begin{scope}[xshift=6.6cm]
    \baIKArm{ba gesucht}{=\,?}{}{}
    \node[ba teiltitel, align=center] at (2.0,-0.6) {b) Inverse Kinematik};
  \end{scope}
\end{tikzpicture}
 in einer figure-Umgebung.
\begin{tikzpicture}[ba]
  % #1 = Stil der Gelenkwinkel, #2 = Zusatz an den Winkeln (z.B. "=?"), #3 = Stil Endeffektor-Label, #4 = Zusatz am Endeffektor
  \newcommand{\baIKArm}[4]{%
    \coordinate (B)  at (0,0);
    \coordinate (J1) at (20:2.3);
    \coordinate (J2) at ($(J1)+(50:2.0)$);
    \coordinate (E)  at ($(J2)+(78:1.8)$);
    % Bezugslinien: Horizontale an der Basis, Verlängerung der Glieder
    \coordinate (H)  at (2.4,0);
    \coordinate (X1) at ($(J1)+(20:1.2)$);
    \coordinate (X2) at ($(J2)+(50:1.1)$);
    \draw[ba hilfslinie, #1] (B) -- (H);
    \draw[ba hilfslinie, #1] (J1) -- (X1);
    \draw[ba hilfslinie, #1] (J2) -- (X2);
    % Glieder
    \draw[ba glied] (B) -- (J1) -- (J2) -- (E);
    % Gelenkwinkel
    \pic[ba winkel, #1, "$\theta_1#2$", angle eccentricity=1.12, pic text options={anchor=west}] {angle = H--B--J1};
    \pic[ba winkel, #1, "$\theta_2#2$", angle eccentricity=1.12, pic text options={anchor=west}] {angle = X1--J1--J2};
    \pic[ba winkel, #1, "$\theta_3#2$", angle eccentricity=1.12, pic text options={anchor=south west}] {angle = X2--J2--E};
    % Gelenke, Basis, Endeffektor
    \node[ba basis]  at (B)  {};
    \node[ba gelenk] at (J1) {};
    \node[ba gelenk] at (J2) {};
    \node[ba endeffektor] at (E) {};
    \node[right=1.5mm of E, #3] {$\vec{r}_\mathrm{E}#4$};
  }
  % a) Vorwärtskinematik: Winkel gegeben, Endeffektorlage gesucht
  \begin{scope}
    \baIKArm{}{}{ba gesucht}{=\,?}
    \node[ba teiltitel, align=center] at (2.0,-0.6) {a) Vorwärtskinematik};
  \end{scope}
  % Trennlinie
  \draw[ba trenner] (5.4,-1.0) -- (5.4,4.9);
  % b) Inverse Kinematik: Endeffektorlage gegeben, Winkel gesucht
  \begin{scope}[xshift=6.6cm]
    \baIKArm{ba gesucht}{=\,?}{}{}
    \node[ba teiltitel, align=center] at (2.0,-0.6) {b) Inverse Kinematik};
  \end{scope}
\end{tikzpicture}
